The architecture's central claim, that appending CST coefficients to
\texttt{mfoil}'s Newton system yields a determined root-find rather than a
search, is falsifiable by a single self-consistency experiment with a
\emph{known} answer, guaranteed realisable by construction. The protocol fits
CST to a known airfoil, runs \texttt{mfoil} in direct mode to generate a
target pressure and edge-velocity distribution from that fit, perturbs the
coefficients or uses the pre-solve of Section~\ref{sec:presolve} to obtain a
starting guess, and asks whether the monolithic inverse Newton iteration
recovers the original coefficient vector $A^\ast$, following the procedure
stated as pseudocode in Section~\ref{sec:algorithms}. Any failure of this test
isolates cleanly to one of the Table~\ref{tab:fm} failure modes rather than to
target infeasibility. This is why this test, and not a search over
independently chosen targets, is the load-bearing evidence for the hypothesis.
It is also, by the same construction, the test's principal weakness, and
Section~\ref{sec:validity} states that weakness in full rather than leaving it
implicit here.

\paragraph{Reference configuration.} NACA 2412, $n=8$ per side, leading edge
prescribed, meaning the first 5 percent of chord is fixed geometrically rather
than left inverse. That is the mixed direct-inverse mitigation for the
leading-edge failure mode of Section~\ref{sec:fm}, and it is the cheapest test
of the whole hypothesis because it isolates the leading-edge question from the
architecture question. Transition is forced during the inverse iterations
(Section~\ref{sec:transition}), target stations are chosen by QR column
pivoting (Section~\ref{sec:ablations}), the iteration is initialised from the
linear pre-solve, and the angle of attack is held fixed. This configuration is
referred to throughout the rest of the paper as the reference configuration.

\paragraph{Protocol, including release-and-verify.} The success criterion is
$\norm{A-A^\ast}_\infty < 10^{-4}$ within nine Newton iterations, together with
a residual history showing sustained superlinear decrease. Because the inverse
solve runs under the forced-transition closure of
Section~\ref{sec:transition}, the protocol additionally requires releasing
transition and re-solving in fully natural, direct mode at the recovered
geometry, then comparing the resulting aerodynamic coefficients against those
of the \emph{reference section} solved in the same natural mode. This is a
check that the forced-transition inverse solve recovered the geometry, rather
than merely finding a zero of the forced-mode residual that happens to be
aerodynamically different.

\paragraph{Result.} Measured against the $10^{-4}$ and nine-iteration
criterion, the solve returns $\norm{A-A^\ast}_\infty = 2.747\x10^{-11}$ in six
Newton iterations, with the target-block residual history
$5.85\x10^{-4} \to 2.05\x10^{-5} \to 1.57\x10^{-7} \to 1.21\x10^{-8} \to
2.13\x10^{-10} \to 9.47\x10^{-12}$, plotted alongside the ablation runs in
Figure~\ref{fig:convergence}. Release-and-verify against a
$\Delta c_l<10^{-3}$, $\Delta c_d<2\x10^{-4}$ criterion gives
$\Delta c_l = 3.4\x10^{-12}$ and $\Delta c_d = 3.0\x10^{-14}$, both essentially
at floating-point noise: the released solve at the recovered geometry
reproduces the reference section's natural-transition coefficients
($c_l=0.449091$, $c_d=0.005774$) to machine precision. Because both sides of
that comparison are natural-mode solves at geometries agreeing to
$2.7\x10^{-11}$, agreement at this level is what smoothness alone would
predict; the check confirms that no aerodynamically distinct root was reached,
and should not be read as independent evidence of accuracy. Realisability,
computed inviscid-consistently as Section~\ref{sec:presolve} requires, is
$3.90\x10^{-4}$; the model gap, diagnostic only, is $0.110$. Sixteen
coefficients are \emph{recovered} by the Newton iteration; the two
leading-edge coefficients $A_{u,0}$ and $A_{l,0}$ are \emph{prescribed} by the
configuration and are not recovered, so all error norms quoted here are over
the 16 free coefficients. Table~\ref{tab:t7} collects the criteria and the
measured values.

\begin{table}[htbp]
\centering
\caption{The falsifiable test at the reference configuration: every criterion
met, each by several orders of magnitude.}
\label{tab:t7}
\begin{tabularx}{\textwidth}{@{}L c c@{}}
\toprule
Criterion & Threshold & Measured \\
\midrule
$\norm{A-A^\ast}_\infty$ over the 16 free coefficients & $<10^{-4}$ & $2.747\x10^{-11}$ \\
Newton iterations & $\leq9$ & $6$ \\
Sustained decrease of the target block & required & 7 decades over 5 steps \\
Release-and-verify $\Delta c_l$ & $<10^{-3}$ & $3.4\x10^{-12}$ \\
Release-and-verify $\Delta c_d$ & $<2\x10^{-4}$ & $3.0\x10^{-14}$ \\
\bottomrule
\end{tabularx}
\end{table}

An independently-run six-vector adversarial review probed circularity,
station-index correspondence, the pinned-coefficient framing, the absence at
the time of a release-and-verify step, realisability semantics, and the
conflation of submap conditioning with system conditioning. Every confirmed
finding is reflected in this paper's reporting discipline: the split between
free and prescribed coefficients above, and the two-metric realisability rule
of Section~\ref{sec:presolve}. The station-index correspondence finding was
closed differently from the others. It was originally answered with a numeric
guard, asserting that the selection-time node and the target's own paneling
agreed to $\max|\Delta x/c|<2\x10^{-5}$. That guard has since been removed,
because addressing a station by surface and chordwise fraction, and
interpolating the target's own pressure curve at that fraction, makes the
correspondence exact rather than approximate, so there is no longer a quantity
for the guard to bound. Two consequences are recorded here rather than left
implicit. Interpolated addressing is what introduces the neglected
$\partial T/\partial A$ term of Section~\ref{sec:newton-system}, and with it
the inexact-Newton behaviour analysed in Section~\ref{sec:h3}. And the runs of
Section~\ref{sec:results}, which predate the change, retain the earlier
nearest-node addressing with a residual correspondence error of up to
$2.0\x10^{-4}$ in chord fraction, ten times the bound the retired guard
imposed; Section~\ref{sec:results} states where that matters.

\begin{figure}[htbp]
\centering
\includegraphics[width=0.75\textwidth]{h3_convergence_overlay.png}
\caption{Newton residual history of the node-indexed baseline run, overlaid
with the evenly-spaced-station wrong-root case and the $n=12$
conditioning-cliff case of Section~\ref{sec:ablations}.}
\label{fig:convergence-overlay}
\end{figure}

Figure~\ref{fig:convergence-overlay} places that history in context by
overlaying the node-indexed baseline against the two ablation runs that fail
for different reasons, the evenly-spaced-station case that converges to a
wrong root and the $n=12$ case that converges slowly; all three are analysed
in Section~\ref{sec:ablations}. A residual history is in any case a scalar
summary. Figures~\ref{fig:t7-geometry} to~\ref{fig:t7-flow} give the geometry
and flow-field view of the same result, regenerated directly from a fresh
re-run of the reference configuration rather than hand-plotted.
Figure~\ref{fig:t7-geometry} overlays the target and Newton-recovered airfoil
outlines, visually indistinguishable at full-chord scale, with a leading-edge
inset showing the same at the scale where the leading-edge failure mode would
be visible if the prescribed-nose mitigation were absent. It also includes a
per-coefficient recovery-error bar chart making explicit which two of eighteen
coefficients are pinned rather than recovered.

\begin{figure}[htbp]
\centering
\includegraphics[width=0.85\textwidth]{paper/fig_t7_geometry_overlay.png}
\caption{Target against Newton-recovered geometry at the reference
configuration, with a leading-edge inset and a per-coefficient recovery-error
bar chart. Grey bars are prescribed, not recovered.}
\label{fig:t7-geometry}
\end{figure}

Figure~\ref{fig:t7-cp} compares target and recovered pressure distributions
both as a full curve and restricted to the 16 QR-pivoted stations the Newton
system actually solved against. The distinction matters: agreement at the
solved stations is enforced by construction, whereas agreement away from them
is not, and it is the full curve that shows the recovered geometry reproduces
the target everywhere rather than only where it was sampled. This is the
pressure-level counterpart of the coefficient-level agreement in
Table~\ref{tab:t7}, and rules out the possibility that the reported error is
an artifact of comparing only at the fitted stations.

\begin{figure}[htbp]
\centering
\includegraphics[width=0.95\textwidth]{paper/fig_t7_cp_comparison.png}
\caption{Target against recovered pressure distribution under forced
transition, plotted in the aerodynamic convention with inverted vertical axis:
full curve (left) and the 16 QR-pivoted target stations only (right).}
\label{fig:t7-cp}
\end{figure}

Figure~\ref{fig:t7-flow} shows the converged, released, natural-transition
flow solution at the recovered geometry: the pressure distribution and the
boundary-layer displacement thickness. This is the physical state from which
the release-and-verify comparison of Table~\ref{tab:t7} was computed, and the
displacement-thickness panel is worth reading alongside
Figure~\ref{fig:viscous}: with transition released, the laminar run extends
well aft of the 5 percent trip used during the inverse iterations, which is
precisely the difference the release step exists to check the solve did not
absorb into geometry.

\begin{figure}[htbp]
\centering
\includegraphics[width=0.95\textwidth]{paper/fig_t7_flow_solution.png}
\caption{Converged flow solution at the recovered geometry after transition is
released: pressure distribution (left) and displacement thickness
$\delta^*/c$ (right).}
\label{fig:t7-flow}
\end{figure}
